% A further Pell shift, independent of the direct quadratic mask improvement.
\section{Shifting the first Pell parameter}\label{sec:shifted-pell}

One more operation can be removed by shifting the first Pell parameter.
Use the notation of Section~\ref{sec:further-pell}, and put
\[
 X=2UM^2=2ws^2r^2n^6,\qquad P=X+1.
\]
Delete the unknown $p$ and its defining equation (E8), and replace
(E9) and (E11) by
\begin{equation}\label{eq:shifted-pell-system}
 X(X+2)k^2+1=\tau^2,\qquad k=r+1+hX.
\end{equation}
Every other equation is retained, including the positive interval
$c=ksn^2+\eta$, $k=\eta+\zeta$.
The mathematical parameter $P$ need not be computed: its coefficient
is $P^2-1=X(X+2)$, and its index modulus is $P-1=X$.

\begin{lemma}\label{lem:shifted-pell}
Under the initial encoding bounds
\[
 3<3L\le B<q\le N\le R,\qquad R,N\ge8,\qquad b\le N,
\]
with $B$ even and $N=q^8$, the replacement
\eqref{eq:shifted-pell-system} preserves existential equivalence with
positive witnesses. It re-chooses only $k,\tau,h,\eta,\zeta$ and removes
or restores $p$. It applies both to $B=Hb^4$, $L=5^{64}$ and to the
direct quadratic encoding $B=Hb^2$, $L=5^{16}$.
\end{lemma}
\begin{proof}
The initial bounds give $U,Y\ge64$, $M=RY\ge512$, and
$A=M(U+1)\ge33280$. Write $C=c$, $K=k$, $J=2R+1$ as before.
The positive interval says $0<C/K-Y<1$. Equation
\eqref{eq:shifted-pell-system} gives $K\ge R+2$, so
$C>(Y-1)K>J$. Proposition~\ref{prop:signed-pell-new} therefore gives
$C=\psi_A(J)$ and $d=\chi_A(J)$, before any approximation is used.
The other Pell equation and its elementary index congruence give
\[
 K=\psi_P(R+1+vX),\qquad v\ge0.
\]
If $v\ge1$, the Pell growth bounds imply
\[
 \frac CK\le\frac{(2A)^{2R}}{(2P-1)^{R+X}}
 \le\left(\frac{2A}{2P-1}\right)^{2R}<\frac12,
\]
contrary to $C/K>Y-1\ge63$. Here $X\ge R$ and
$2A/(2P-1)=2M(U+1)/(4UM^2+1)<1/2$. Thus
$K=\psi_P(R+1)$.

Put $\xi=(U+1)^{2R}/U^R$. At these exact indices the same growth bounds
give
\begin{align}
 \frac CK
 &\le\frac{(2A)^{2R}}{(2P-1)^R}
   =\xi\left(1+\frac1{4UM^2}\right)^{-R}<\xi,
       \label{eq:shifted-upper}\\
 \frac CK
 &\ge\xi\left(1-\frac1{2A}\right)^{2R}
             \left(1+\frac1X\right)^{-R}
  >\xi\left(1-\frac RA-\frac RX\right).
       \label{eq:shifted-lower}
\end{align}
The last inequality follows from Bernoulli's inequality, with both
factors positive. Since $A/X=(U+1)/(2UM)<1$, the error coefficient is
\[
 \varepsilon=\frac RA+\frac RX
 <\frac{2}{Y(U+1)}<\frac12.
\]
Consequently $C/K>\xi/2>U^R/2>U^{R-1}+1$, while $C/K<Y+1$.
Hence $Y>U^{R-1}$ and $A>RU^R$.

This last bound permits the two exponent arguments in their original
order. First,
\[
 2^{3J}\le8N^{2R}\le RU^R<A,\qquad (bw)^3\le U^R<A,
\]
so (E14) and the $\chi$ form of Lemma~2.22 of~\cite{J82} give
$bw=2^J$, hence
\begin{equation}\label{eq:shifted-u-bound}
 U\ge Nbw\ge8\cdot2^J>4\cdot2^{2R}.
\end{equation}
Next, $\kappa<C$, (E19), (E20), and $0<L<J<A$ give
$\kappa=\psi_A(L)$. The bounds
$B^{3L},q^3\le N^N\le U^R<A$ and (E18) give $q=B^L$ by the same
exponent lemma. Since $B$ is even already, $q$ and $N$ are now even.
Neither argument has used parity rounding or central-binomial
divisibility.

From $\xi/2<C/K<Y+1$ we have $\xi<2Y+2<3Y$, so
\eqref{eq:shifted-upper}--\eqref{eq:shifted-lower} imply
\[
 0<\xi-C/K<\xi\varepsilon
 <\frac6{U+1}\le\frac6{65}<\frac14.
\]
The binomial expansion, using \eqref{eq:shifted-u-bound}, gives
\[
 0<\xi-\lfloor\xi\rfloor<\frac14,
 \qquad \lfloor\xi\rfloor\equiv\binom{2R}{R}\pmod U.
\]
Thus $|Y-\lfloor\xi\rfloor|<3/2<2$. Both integers are even, since
$Y,U$ are even and $\binom{2R}{R}$ is even. Therefore
$Y=\lfloor\xi\rfloor$, so $N^2\mid\binom{2R}{R}$ as required.

It remains to justify strict positivity when changing witnesses.
Start with an old solution and set
\[
 K_{\rm new}=\psi_{X+1}(R+1),\quad
 \tau_{\rm new}=\chi_{X+1}(R+1),\quad
 h_{\rm new}=\frac{K_{\rm new}-R-1}{X}>0.
\]
The quotient is integral by the Pell congruence and positive by strict
growth. The old solution has $Y=\lfloor\xi\rfloor$ and $\xi<2Y$.
Thus the new ratio estimates give
$0<\xi-C/K_{\rm new}<4/U$, whereas the positive fractional tail gives
\[
 \xi-Y\ge\frac{\binom{2R}{R-1}}U\ge\frac{2R}U>\frac4U.
\]
Consequently $Y<C/K_{\rm new}<\xi<Y+1/4$, and
$\eta_{\rm new}=C-K_{\rm new}Y$,
$\zeta_{\rm new}=K_{\rm new}-\eta_{\rm new}$ are positive integers.

Conversely, restore $p=X$ and choose
$K_{\rm old}=\psi_X(R+1)$, $\tau_{\rm old}=\chi_X(R+1)$, and
$h_{\rm old}=(K_{\rm old}-R-1)/(X-1)>0$. The elementary old-base
estimates at these exact indices give
\[
 \xi(1-R/A)<C/K_{\rm old}<\xi(1+R/X),
 \qquad |C/K_{\rm old}-\xi|<2/U.
\]
The same positive tail yields $C/K_{\rm old}>Y$, while its upper
bound gives $C/K_{\rm old}-Y<1/4+2/U<1$. The corresponding old
$\eta,\zeta$ are therefore positive as well. Every other witness is
preserved. The proof uses only the stated initial bounds and evenness
of $B$, which were established for both encodings independently of
the Pell block.
\end{proof}

\begin{theorem}\label{thm:best112}
For each r.e.\ set, the effective admissible index $(Z,V,H)$ of
Theorem~\ref{thm:best113} gives a system in 32 positive unknowns and
20 equations such that, for positive input $x$, membership is equivalent
to a valid straight-line certificate with \textbf{112 arithmetic
instructions}: 62 multiplications and 50 additions. Equality tests and
supplied parameters are free. Constructing $2,4,5,L=5^{16}$ from $1$
adds seven instructions, giving 119 under that convention.
\end{theorem}
\begin{proof}
Apply Lemma~\ref{lem:shifted-pell} to the 113-operation direct quadratic
system of Theorem~\ref{thm:best113}. The shared value $X$ was already
computed by the old (E8) block. Replace the two instructions for
$p^2-1$ by
\[
 X_+=X+2,\qquad D_P=XX_+,
\]
and use $hX$ in place of $h(p-1)$. This deletes one subtraction,
the positive unknown $p$, and its free defining equality. All other
counts are unchanged. The complete primitive schedule, equality list,
declared domains, and exact residual checks are in
\file{round6\_1980\_certificate.py} and its JSON receipt. The standalone
shift and its full positive-witness proof are also recorded in
\file{round5\_1980\_shifted\_pell\_certificate.py} and
\file{PELL\_SHIFTED\_BASE\_PROOF.md}.
\end{proof}
